% ===== PRÉAMBULE CANONIQUE — à coller VERBATIM en tête de CHAQUE .tex =====
\documentclass[11pt,a4paper]{article}
\usepackage[utf8]{inputenc}\usepackage[T1]{fontenc}\usepackage{lmodern}
\usepackage[french]{babel}
\usepackage{amsmath,amssymb,mathtools}\usepackage{icomma}
\usepackage[version=4]{mhchem}          % \ce{...} : équations & formules
\usepackage{siunitx}\sisetup{locale=FR} % \qty{}{}, \si{}, \ang{}
\usepackage{chemfig}                    % \chemfig{...} : structures développées
\usepackage{xcolor}\usepackage{colortbl}
\usepackage{tikz}\usepackage{pgfplots}\pgfplotsset{compat=1.18}
\usepackage[most]{tcolorbox}\usepackage{enumitem}\usepackage{array,booktabs}
\usepackage[a4paper,margin=2.2cm]{geometry}
\usetikzlibrary{arrows.meta,calc,angles,quotes,positioning,patterns,backgrounds,shapes.geometric}
\definecolor{primary}{RGB}{222,49,99}\definecolor{primarydark}{RGB}{150,22,55}
\definecolor{defcol}{RGB}{0,114,178}\definecolor{methcol}{RGB}{52,92,148}
\definecolor{excol}{RGB}{0,135,90}\definecolor{attncol}{RGB}{200,40,55}
\tcbset{boxbase/.style={enhanced,breakable,boxrule=0.9pt,arc=2.5pt,
  left=3mm,right=3mm,top=2.3mm,bottom=2.3mm,fonttitle=\bfseries,coltitle=white}}
% --- boîtes TUTORIEL (noms EXACTS, ne pas renommer) ---
\newtcolorbox{cle}[1][]{boxbase,colback=defcol!6,colframe=defcol,
  title={\textsc{À retenir}\ifx\relax#1\relax\relax\else\textmd{ : \itshape #1}\fi}}
\newtcolorbox{methode}[1][]{boxbase,colback=methcol!7,colframe=methcol,sharp corners,
  title={\textsc{Méthode}\ifx\relax#1\relax\relax\else\textmd{ --- \itshape #1}\fi}}
\newtcolorbox{exemplebox}[1][]{boxbase,colback=excol!5,colframe=excol,
  title={\textsc{Exemple}\ifx\relax#1\relax\relax\else\textmd{ : \itshape #1}\fi}}
\newtcolorbox{piege}[1][]{boxbase,colback=attncol!6,colframe=attncol,
  title={\textsc{Piège}\ifx\relax#1\relax\relax\else\textmd{ : \itshape #1}\fi}}
\newtcolorbox{remarque}[1][]{boxbase,colback=black!4,colframe=black!45,
  title={\textsc{Remarque}\ifx\relax#1\relax\relax\else\textmd{ : \itshape #1}\fi}}
% --- boîtes EXERCICE / CORRIGÉ (noms EXACTS, auto-comptées) ---
\newtcolorbox[auto counter]{exo}[1][]{boxbase,colback=primary!4,colframe=primary,
  title={\textsc{Exercice}~\thetcbcounter\ifx\relax#1\relax\relax\else\textmd{ --- \itshape #1}\fi}}
\newtcolorbox[auto counter]{corrige}[1][]{boxbase,colback=excol!4,colframe=excol,
  title={\textsc{Corrigé}~\thetcbcounter\ifx\relax#1\relax\relax\else\textmd{ --- \itshape #1}\fi}}
\newcommand{\R}{\mathbb{R}}\newcommand{\N}{\mathbb{N}}\newcommand{\Z}{\mathbb{Z}}\newcommand{\ds}{\displaystyle}
% =========================================================================
\begin{document}
\begin{corrige}[constante à partir d'un taux de décomposition]
$\qty{40}{\percent}$ de $\qty{0.50}{\mol}$ décomposé $=\qty{0.20}{\mol}$.
\begin{enumerate}[label=\alph*)]
  \item tableau d'avancement (en $\si{\mol}$) :
  \begin{center}
  \begin{tabular}{l|ccc}
   & $\ce{PCl5}$ & $\ce{PCl3}$ & $\ce{Cl2}$\\\hline
  initial     & $0{,}50$ & $0$ & $0$\\
  transformé  & $-0{,}20$ & $+0{,}20$ & $+0{,}20$\\
  équilibre   & $0{,}30$ & $0{,}20$ & $0{,}20$\\
  \end{tabular}
  \end{center}
  \item dans $\qty{1.0}{\litre}$ : $[\ce{PCl5}]=\qty{0.30}{\mol\per\litre}$, $[\ce{PCl3}]=[\ce{Cl2}]=\qty{0.20}{\mol\per\litre}$.
  \item $K_c = \dfrac{[\ce{PCl3}]\,[\ce{Cl2}]}{[\ce{PCl5}]} = \dfrac{0{,}20\times 0{,}20}{0{,}30} = \dfrac{0{,}040}{0{,}30} \approx 0{,}13$.
\end{enumerate}
\end{corrige}

\begin{corrige}[concentration à l'équilibre : une équation du second degré]
\begin{enumerate}[label=\alph*)]
  \item dans $\qty{1.0}{\litre}$ : $[\ce{N2O4}] = 1{,}0-x$ et $[\ce{NO2}] = 2x$.
  \item $K = \dfrac{[\ce{NO2}]^{2}}{[\ce{N2O4}]} = \dfrac{(2x)^{2}}{1-x} = \dfrac{4x^{2}}{1-x} = 2{,}0$, d'où $4x^{2} = 2(1-x)$, soit
  \[ 2x^{2} + x - 1 = 0 \;\Rightarrow\; x = \frac{-1 + \sqrt{1 + 8}}{4} = \frac{-1+3}{4} = 0{,}50 \quad(\text{la racine négative est rejetée}). \]
  \item $[\ce{NO2}] = 2x = \qty{1.0}{\mol\per\litre}$ et $[\ce{N2O4}] = 1-x = \qty{0.50}{\mol\per\litre}$ (vérification : $1{,}0^{2}/0{,}50 = 2{,}0$).
\end{enumerate}
\end{corrige}

\begin{corrige}[équilibre hétérogène : pression totale]
\begin{enumerate}[label=\alph*)]
  \item les solides n'apparaissent pas : $K_p = p_{\ce{H2O}}\times p_{\ce{CO2}}$.
  \item partant du solide seul, $\ce{H2O}$ et $\ce{CO2}$ sont produits en quantités \emph{égales} (1 mol chacun) : $p_{\ce{H2O}} = p_{\ce{CO2}} = p$.
  \item $K_p = p^{2} = 0{,}25 \Rightarrow p = \qty{0.5}{atm}$. Pression totale $= p_{\ce{H2O}} + p_{\ce{CO2}} = 0{,}5 + 0{,}5 = \qty{1.0}{atm}$.
\end{enumerate}
\end{corrige}

\begin{corrige}[Le Chatelier : réaction endothermique]
Gauche : $1$ mol de gaz ($\ce{CO2}$) ; droite : $2$ mol de gaz ($\ce{CO}$).
\begin{enumerate}[label=\alph*)]
  \item \textbf{Oui} : la réaction est endothermique, $T\nearrow$ la déplace vers la droite.
  \item \textbf{Non} : un catalyseur ne déplace pas l'équilibre.
  \item \textbf{Oui} : augmenter le volume abaisse la pression, ce qui favorise le côté à \emph{plus} de gaz, soit la droite.
  \item \textbf{Non} : le carbone est un solide pur, il n'intervient pas.
  \item \textbf{Oui} : retirer un produit déplace l'équilibre vers la droite pour le reformer.
\end{enumerate}
\end{corrige}

\begin{corrige}[Le Chatelier : réaction exothermique (procédé de contact)]
Gauche : $3$ mol de gaz ; droite : $2$ mol de gaz.
\begin{enumerate}[label=\alph*)]
  \item \textbf{Oui} : ajouter un réactif ($\ce{O2}$) déplace l'équilibre vers la droite (il est consommé).
  \item \textbf{Oui} : la réaction est exothermique, abaisser $T$ la déplace vers la droite.
  \item \textbf{Oui} : $p\nearrow$ favorise le côté à moins de gaz, soit la droite ($3\to 2$).
  \item \textbf{Non} : un catalyseur accélère l'équilibre sans le déplacer (aucun effet sur le rendement).
  \item \textbf{Non} : à volume constant, l'argon ne change pas les pressions partielles ; aucun déplacement.
\end{enumerate}
\end{corrige}

\end{document}
